% !PRINTED 572
\weberpage{572}
That these forms are, however, really different from one another
follows from the fact that (by \S.\,145, 150) we can put for the
variables $t_1, t_2, \ldots, t_n$ integral rational numbers such
that $\tau$ passes over into a primitive root $\gamma$ of the prime
ideal $\mathfrak p$. By the same substitution the quantities (10)
become
\weq{11}{\equiv \gamma, \gamma^p, \gamma^{p^2}, \ldots, \gamma^{p^{f-1}}\ (\mathrm{mod}\ \mathfrak p),}
which are all different from one another with respect to the modulus
$\mathfrak p$. Hence, too, no two of the forms (10) can be congruent
with respect to the modulus $\mathfrak p$, since otherwise the two
corresponding numbers (11) would also turn out congruent.

We sum this up as a theorem thus:

1. \gesperrt{Every congruence (4), of which one root is $t = \tau$,
has the $f$ different roots}
\[
\tau, \tau_1, \tau_2, \ldots, \tau_{f-1}.
\]

Accordingly, dividing $\Phi(t)$ algebraically by $t - \tau$, we can
set
\[
\Phi(t) \equiv (t - \tau)\, \Phi_1(t)\ (\mathrm{mod}\ \mathfrak p)
\]
and $\tau_1, \tau_2, \ldots, \tau_{f-1}$ are roots of
$\Phi_1(t) \equiv 0$, in which, however, $\Phi_1(t)$ does not yet
have rational coefficients, but coefficients in $\mathfrak o$. If we
divide $\Phi_1(t)$ again by $t - \tau_1$, and continue in this way,
it follows finally that, if we now set the function of degree $f$
\weq{12}{\Pi(t) = (t - \tau)(t - \tau_1) \ldots (t - \tau_{f-1})}
then
\weq{13}{\Phi(t) \equiv \Pi(t)\, \Phi_0(t)\ (\mathrm{mod}\ \mathfrak p),}
in which $\Phi_0(t)$ is an integral function with coefficients in
$\mathfrak o$.

The function $\Pi(t)$ depends on the variables $t, t_1, \ldots, t_n$,
and to express this we set
\[
\Pi(t) = \Pi(t, t_1, \ldots, t_n).
\]

The coefficients of this form are numbers in $\mathfrak o$, and it
can further be shown that they are congruent with integral rational
numbers with respect to the modulus $\mathfrak p$. This proof is
obtained by raising to the $p$th power:
\[
[\Pi(t)]^p \equiv (t^p - \tau^p)(t^p - \tau_1^p) \ldots (t^p - \tau_{f-1}^p).
\]

Now, however, by (9)
\[
\tau_r^p \equiv \omega_1^{p^{r+1}} t_1^p + \omega_2^{p^{r+1}} t_2^p + \cdots + \omega_n^{p^{r+1}} t_n^p\ (\mathrm{mod}\ \mathfrak p),
\]
and so we obtain $\tau_r^p$ from $\tau_{r+1}$ when we replace $t_i$
by $t_i^p$,
%CONT>>
